% GENERATED FILE. Edit canonical lessons, then run npm run generate:books.
% canonical-source-hash: fnv1a64:9e1dd6581f5a957e
% canonical-lessons: TA-C216-inaiyam, TA-C216-katavuccol, TA-C216-nalital, TA-C216-cittu, TA-C216-varaipatam

\chapter{Internet, Password, Newspaper, Ticket, Map}
\label{ch:ta-internet-password-newspaper-ticket-map}
\hlchaptermodality{216}

\begin{chapteropening}
\textbf{By the end of this chapter:} \emph{I can say the internet, a password, a newspaper, a ticket and a map.}

Complete the last lesson of chapter 216: I can say the internet, a password, a newspaper, a ticket and a map.
\end{chapteropening}

\section[\texorpdfstring{\ta{இணையம்} (iṇaiyam)}{இணையம் (iṇaiyam)}]{\ta{இணையம்} (iṇaiyam) — the internet}

\label{lesson:TA-C216-inaiyam}



\begin{quote}
Before the new one: say the Tamil for a letter, then the Tamil for an envelope; a cover.
\end{quote}

\subsection*{You'll want to know: \ta{இணையம்}}
\textbf{\ta{இணையம்}} — \emph{iṇaiyam} — \textquotedblleft{}the internet\textquotedblright{}.

\begin{grammarlens}[title={What you've built}]
One more word for this chapter.
\end{grammarlens}

\subsection*{Your turn}
\begin{itemize}
\raggedright
  \item \emph{Say it:} \emph{iṇaiyam}
  \item \emph{Say it:} \emph{iṇaiyam}, once more
  \item \emph{Say it:} say \emph{iṇaiyam}
  \item \emph{Recall:} say the Tamil for a letter, then the Tamil for an envelope; a cover, then say \emph{iṇaiyam} again
\end{itemize}

\begin{tcolorbox}[breakable,colback=teal!4,colframe=teal!35!black,title={Before you move on}]
What does \ta{இணையம்} mean? (\textquotedblleft{}The internet\textquotedblright{}.) Say it once more.
\end{tcolorbox}
\section[\texorpdfstring{\ta{கடவுச்சொல்} (kaṭavuccol)}{கடவுச்சொல் (kaṭavuccol)}]{\ta{கடவுச்சொல்} (kaṭavuccol) — a password}

\label{lesson:TA-C216-katavuccol}



\begin{quote}
Before the new one: say the Tamil for an envelope; a cover, then the Tamil for the internet.
\end{quote}

\subsection*{You'll want to know: \ta{கடவுச்சொல்}}
\textbf{\ta{கடவுச்சொல்}} — \emph{kaṭavuccol} — \textquotedblleft{}a password\textquotedblright{}.

\begin{grammarlens}[title={What you've built}]
One more word for this chapter.
\end{grammarlens}

\subsection*{Your turn}
\begin{itemize}
\raggedright
  \item \emph{Say it:} \emph{kaṭavuccol}
  \item \emph{Say it:} \emph{kaṭavuccol}, once more
  \item \emph{Say it:} say \emph{kaṭavuccol}
  \item \emph{Recall:} say the Tamil for an envelope; a cover, then the Tamil for the internet, then say \emph{kaṭavuccol} again
\end{itemize}

\begin{tcolorbox}[breakable,colback=teal!4,colframe=teal!35!black,title={Before you move on}]
What does \ta{கடவுச்சொல்} mean? (\textquotedblleft{}A password\textquotedblright{}.) Say it once more.
\end{tcolorbox}
\section[\texorpdfstring{\ta{நாளிதழ்} (nāḷitaḻ)}{நாளிதழ் (nāḷitaḻ)}]{\ta{நாளிதழ்} (nāḷitaḻ) — a newspaper}

\label{lesson:TA-C216-nalital}



\begin{quote}
Before the new one: say the Tamil for the internet, then the Tamil for a password.
\end{quote}

\subsection*{You'll want to know: \ta{நாளிதழ்}}
\textbf{\ta{நாளிதழ்}} — \emph{nāḷitaḻ} — \textquotedblleft{}a newspaper\textquotedblright{}.

\begin{grammarlens}[title={What you've built}]
One more word for this chapter.
\end{grammarlens}

\subsection*{Your turn}
\begin{itemize}
\raggedright
  \item \emph{Say it:} \emph{nāḷitaḻ}
  \item \emph{Say it:} \emph{nāḷitaḻ}, once more
  \item \emph{Say it:} say \emph{nāḷitaḻ}
  \item \emph{Recall:} say the Tamil for the internet, then the Tamil for a password, then say \emph{nāḷitaḻ} again
\end{itemize}

\begin{tcolorbox}[breakable,colback=teal!4,colframe=teal!35!black,title={Before you move on}]
What does \ta{நாளிதழ்} mean? (\textquotedblleft{}A newspaper\textquotedblright{}.) Say it once more.
\end{tcolorbox}
\section[\texorpdfstring{\ta{சீட்டு} (cīṭṭu)}{சீட்டு (cīṭṭu)}]{\ta{சீட்டு} (cīṭṭu) — a ticket; a slip of paper}

\label{lesson:TA-C216-cittu}



\begin{quote}
Before the new one: say the Tamil for a password, then the Tamil for a newspaper.
\end{quote}

\subsection*{You'll want to know: \ta{சீட்டு}}
\textbf{\ta{சீட்டு}} — \emph{cīṭṭu} — \textquotedblleft{}a ticket; a slip of paper\textquotedblright{}.

\begin{grammarlens}[title={What you've built}]
One more word for this chapter.
\end{grammarlens}

\subsection*{Your turn}
\begin{itemize}
\raggedright
  \item \emph{Say it:} \emph{cīṭṭu}
  \item \emph{Say it:} \emph{cīṭṭu}, once more
  \item \emph{Say it:} say \emph{cīṭṭu}
  \item \emph{Recall:} say the Tamil for a password, then the Tamil for a newspaper, then say \emph{cīṭṭu} again
\end{itemize}

\begin{tcolorbox}[breakable,colback=teal!4,colframe=teal!35!black,title={Before you move on}]
What does \ta{சீட்டு} mean? (\textquotedblleft{}A ticket; a slip of paper\textquotedblright{}.) Say it once more.
\end{tcolorbox}
\section[\texorpdfstring{\ta{வரைபடம்} (varaipaṭam)}{வரைபடம் (varaipaṭam)}]{\ta{வரைபடம்} (varaipaṭam) — a map}

\label{lesson:TA-C216-varaipatam}



\begin{quote}
Before the new one: say the Tamil for a newspaper, then the Tamil for a ticket; a slip of paper.
\end{quote}

\subsection*{You'll want to know: \ta{வரைபடம்}}
\textbf{\ta{வரைபடம்}} — \emph{varaipaṭam} — \textquotedblleft{}a map\textquotedblright{}.

\begin{grammarlens}[title={What you've built}]
One more word for this chapter.
\end{grammarlens}

\subsection*{Your turn}
\begin{itemize}
\raggedright
  \item \emph{Say it:} \emph{varaipaṭam}
  \item \emph{Say it:} \emph{varaipaṭam}, once more
  \item \emph{Say it:} say \emph{varaipaṭam}
  \item \emph{Recall:} say the Tamil for a newspaper, then the Tamil for a ticket; a slip of paper, then say \emph{varaipaṭam} again
\end{itemize}

\begin{tcolorbox}[breakable,colback=teal!4,colframe=teal!35!black,title={Before you move on}]
What does \ta{வரைபடம்} mean? (\textquotedblleft{}A map\textquotedblright{}.) Say it once more.
\end{tcolorbox}
